% Folder 144, batch 09 (archivists' pages 161-180).
% Pass: Opus 5.5 (claude-opus-5-5), 2026-10-03 — first pass, unchecked against the pages by a human.
%
% Skipped, no hand on them: 162, 164, 166, 168, 170, 175, 177, 180 (typed
% university notices on the versos) and 171, 178 (blank separator sheets).
% Page 173 is a typed seminar notice scanned upside down; his notes on it are
% transcribed. The pencilled numbers on the leaves run one ahead (162-181);
% \page{} follows the batch numbering given by the split (page 1 = 161).
\documentclass[11pt,a4paper]{article}
\input{../preamble/grothendieck}

\folder{144}
\batch{9}
\pages{161}{180}
\foldertitle{[Suite autour de Teichmüller dont] Teichmülleries, 1981-1982 : notes manuscrites (1981-1983, s.d.), lettres (1981, s.d.).}
\dating{1981-1983}
\watermark{Édition de démonstration}

\begin{document}

\section{Le groupe $\Gamma_{1,1}$ : générateurs $\tau'$, $\tau''$, $\sigma$, $\rho'$, $\rho''$}

\page{161}
\note{Le feuillet commence au milieu d'un calcul : les objets $X'$, $X''$, $\varepsilon_0$, $\varepsilon_1$, $\omega$ sont introduits avant ce lot.}

\drawing{à gauche, un parallélépipède en pointillés, flèches le long des arêtes ; au centre, un cylindre (anneau) avec deux arcs $X'$ (en dessus) et $X''$ (le long du cylindre), les points $a_0, a_1, a_2, a_3$, deux vecteurs $e'$, $e''$ issus de $a_0$ et une flèche de rotation $\sigma$ ; au-dessus, une petite croix fléchée.}

\[
\begin{array}{ll}
X'_{\xi} & \varepsilon_0(X') \xrightarrow{\ \sim\ } \varepsilon_1(X'') \\
X''_{\xi} & \varepsilon_1(X') \xleftarrow{\ \sim\ } \varepsilon_0(X'') \\
\hline
 & \omega(X') \simeq \omega(X'')
\end{array}
\qquad \Longrightarrow
\]
\note{sous la dernière ligne, une flèche vers $\omega(X' \wedge X'')$, avec un signe $(-)$ entouré et la mention « anticommutatif ».}

$\tau'$, $\tau''$, $\sigma$ ne dépendent pas du « choix d'orientations ».

$e'$, $e''$ fixent les orientations \uncertain{et aussi} une 2-orientation \struck{\ill{}} sur $X'$, $X''$, d'où $\rho'$, $\rho''$.
\[
\begin{aligned}
\tau' e' &= e' & \qquad \tau'' e' &= -e' \\
\tau' e'' &= -e'' & \qquad \tau'' e'' &= e''
\end{aligned}
\]
\note{dans $\tau'' e'' = e''$, un signe biffé devant le second membre.}

d'où $\rho'$, $\rho''$ :
\[
\left\{
\begin{aligned}
&\tau'(\rho') = \tau''(\rho') = \rho'^{-1} \\
&\tau'(\rho'') = \tau''(\rho'') = \rho''^{-1} \\
&\sigma(\rho') = \rho'', \quad \sigma\rho'' = \rho'
\end{aligned}
\right.
\qquad
\begin{aligned}
\rho'(e') &= e' \\
\rho'(e'') &= e' + e''
\end{aligned}
\]

\page{163}
$\tau'$, $\tau''$, $\sigma$, $\rho'$, $\rho''$ \note{$\tau''$ et $\sigma$ soulignés ; un $\sigma$ biffé dessous.}
\[
\begin{aligned}
\rho' &\longmapsto \varepsilon_0 \\
\rho'' &\longmapsto \varepsilon_1 \\
\tau'' &\longmapsto \tau_\infty = \tau = r_\infty \\
\tau' &\longmapsto \omega_0 \tau_\infty = -\tau \\
\sigma &\longmapsto \sigma_\infty = \sigma
\end{aligned}
\qquad
\left\{
\begin{aligned}
&\tau'^2 = \tau''^2 = \sigma^4 = 1, \quad \sigma^2 = \tau'\tau'' \\
&\tau'(\rho') = \tau''(\rho') = \rho'^{-1}, \quad \tau'(\rho'') = \tau''(\rho'') = \rho''^{-1} \\
&\sigma(\rho') = \rho'' \quad (\Rightarrow \sigma(\rho'') = \rho')
\end{aligned}
\right.
\]
\note{à la fin de la première ligne de l'accolade, $\sigma(\rho') = \rho$ biffé.}

$(\sigma\rho')^3 \neq 1$ dans $\Gamma^{+}_{1,1}$ (mais \add{$= 1$} dans un groupe \uncertain{spécial}).
\note{le signe entre $(\sigma\rho')^3$ et $1$ est surchargé ; la lecture $\neq$ est douteuse.}
\[
\begin{aligned}
\sigma^{-1}\rho' = \sigma^3\rho' &\longmapsto \sigma^3\varepsilon_0 = \sigma^3(\sigma\rho) = \rho \\
\tau''\sigma &\longmapsto \tau'_\infty = r_0 \\
\tau''\rho' = \underbrace{(\tau''\sigma)}_{\mapsto r_0}\underbrace{(\sigma^{-1}\rho')}_{\mapsto \rho} &\longmapsto r_1
\end{aligned}
\]
Encadré :
\[
\begin{array}{l|l}
\tau'' \longmapsto r_\infty & \rho' \longmapsto \varepsilon_0 \\
\tau''\sigma \longmapsto r_0 & \sigma \longmapsto \sigma \\
\tau''\rho' \longmapsto r_1 & \sigma^{-1}\rho' \longmapsto \rho
\end{array}
\]

\drawing{à gauche, un tore dessiné en perspective, avec deux cercles méridiens et des courbes en pointillés ; à droite, un carré de sommets $a_0, a_1, a_2, a_3$ (côté $a_0a_1$ fléché), un petit carré intérieur, deux axes fléchés $e'$ (horizontal) et $e''$ (vertical), avec les légendes $\tau' = \mathrm{id}$ et $\tau'' = \ill{}$.}

$\mathfrak{X} = (X, X', X'')$, $\mathcal{E}(\mathfrak{X})$.

$\{X', X''\}$ = ens. des codiagonales du carré

$\varepsilon_0(X') \simeq$ ens. des deux côtés opposés du carré qui correspondent à la $X''$

$\varepsilon_1(X') \simeq$ ens. des deux côtés opposés du carré qui correspondent à $X'$
\[
\begin{aligned}
\varepsilon_0(X') \times \varepsilon_1(X') &\simeq \text{ens. des \struck{côtés} \add{\uncertain{sommets}} du carré} \\
&\simeq \varepsilon_1(X'') \times \varepsilon_0(X'') \\
&\simeq \varepsilon_0(X') \times \varepsilon_0(X'') \\
&\simeq \varepsilon_1(X') \times \varepsilon_1(X'')
\end{aligned}
\]
\note{avant la dernière ligne, un début $\simeq \varepsilon_0(X'') \times \ldots$ biffé. En bas à gauche : $\mathcal{E}(\mathfrak{X})$ suivi de $\varepsilon_0(X')$ biffé, et un mot biffé \ill{}.}

\page{165}
\[
\left\{
\begin{aligned}
&\sigma_1^2 = \sigma_2^2 = (\sigma_1\sigma_2)^2 = 1 \\
&\rho_1, \rho_2
\end{aligned}
\right.
\qquad
\sigma_i\rho_j = \rho_j^{-1}
\]
\[
\tilde\tau\tilde\rho = \tilde\rho^{-1}, \qquad \tilde\tau\tilde\sigma = \tilde\sigma^{-1}, \qquad \tilde\rho^3 -
\]

\drawing{un cylindre ; au-dessous, un empilement de boîtes en perspective avec des axes $x$, $y$ ; à gauche, un cylindre sur lequel est collée une anse, avec une flèche de rotation et deux vecteurs $e_1$, $e_2$ ; plus bas, un parallélépipède hachuré avec les points $a_0, a_1, a_2, a_3$, $b_0, b_1, b_2, b_3$, $c_0, c_1, c_2, c_3$ ; à côté, $\sigma\varphi$ et $D_\sigma$ (un premier mot biffé).}
\[
x \longmapsto -x, \qquad y \longmapsto y, \qquad (x, y, z) \longmapsto ix,
\]
\[
(z, t) \longrightarrow (iz, -t)
\]
Encadré :
\[
\sigma^2 = \sigma_1\sigma_2 = \sigma_2\sigma_1
\]
\[
\sigma_1^2 = \sigma_2^2 \ (= \sigma^4) = 1
\]
\[
\sigma_1(\sigma) = \sigma^{-1}, \qquad \sigma_2(\sigma) = \sigma^{-1} \quad ?
\]
\[
\text{\struck{$(\sigma_1, \sigma_2)$}} \qquad \text{\struck{$(\sigma_1\sigma_2)^2$}}
\]
\[
\sigma_2 = \sigma_1\sigma^2 = \sigma^2\sigma_1
\]
\[
\begin{pmatrix} 1 & 1 \\ 0 & 1 \end{pmatrix}
\begin{pmatrix} 1 & -1 \\ 0 & 1 \end{pmatrix}
\]
\[
\begin{aligned}
\sigma(z, t) &= (iz, -t) \\
\sigma_1(z, t) &= (\bar z, t) \\
\sigma_2(z, t) &= (-\bar z, t)
\end{aligned}
\]
\[
(z, t) \longmapsto (\bar z, t) \longmapsto (i\bar z, -t) \longmapsto (-iz, -t)
\]

\page{167}
\[
\sigma_1, \sigma_2, \sigma, \rho_1, \rho_2
\]
\[
\left.
\begin{aligned}
&\sigma_1^2 = \sigma_2^2 = (\sigma_1\sigma_2)^2 = 1 \\
&\sigma_1\sigma_2 = \sigma^2, \quad \sigma_1(\sigma) = \sigma_2(\sigma) = \sigma^{-1}
\end{aligned}
\right\}
\quad D_\sigma \simeq D_4
\]
$D_\sigma$ engendré par $\sigma_1$, $\sigma_2$, $\sigma$.
\[
\left.
\begin{aligned}
&\sigma_i(\rho_j) = \rho_j^{-1} \\
&\sigma(\rho_1) = \rho_2
\end{aligned}
\right|
\Longrightarrow \sigma(\rho_2) = \rho_1
\]
$\sigma$ : $D_\sigma \ast_{D_{\sigma_1}} D_{\rho_1}$ ($D_{\rho_1}$ engendré par $\rho_1$, $\sigma_1$)

Encadré :
\[
\sigma, \sigma_1, \rho_1 \;\Big|\; \sigma^4 = \sigma_1^2 = 1, \qquad
\begin{aligned}
\sigma_1(\sigma) &= \sigma^{-1} \\
\sigma_1(\rho_1) &= \rho_1^{-1} \\
\sigma^2(\rho_1) &= \rho_1
\end{aligned}
\]
$\sigma_2 = \sigma_1\sigma^2$ ($\sigma^2$ \uncertain{central}) \uncertain{commute} à $\sigma_1$ ; $\sigma_1\sigma_2 =$

$\sigma_2(\rho_1) =$

\drawing{$\sigma^2$ — un cercle traversé par deux diamètres perpendiculaires ; plus bas, un hexagone dont les côtés ne se rejoignent pas, et un cercle portant les points $a_1, c_0, b_0, a_0, c_3, b_3, a_3, c_2, b_2, a_2, c_1, b_1$, traversé de deux droites pleines et de segments en pointillés.}
\[
\left\{
\begin{aligned}
\sigma_1 e_1 &= e_1 \\
e_1 e_2 &= -e_2
\end{aligned}
\right.
\qquad
\left\{
\begin{aligned}
\sigma_2 e_1 &= -e_1 \\
\sigma_2 e_2 &= e_2
\end{aligned}
\right.
\qquad
\left\{
\begin{aligned}
\sigma e_1 &= e_2 \\
\sigma e_2 &= -e_1
\end{aligned}
\right.
\]
\note{« $e_1 e_2 = -e_2$ » tel quel ; on attend $\sigma_1 e_2 = -e_2$.}
\[
\rho_1 e_1 = e_1, \qquad \rho_1 e_2 = e_2 - e_1
\]
\note{dans la dernière formule, un premier second membre biffé.}
\[
\tau' = \begin{pmatrix} 1 & 0 \\ 0 & -1 \end{pmatrix}, \qquad
\tau'' = \begin{pmatrix} -1 & 0 \\ 0 & 1 \end{pmatrix} = \tau_\infty, \qquad
\sigma = \begin{pmatrix} 0 & -1 \\ 1 & 0 \end{pmatrix}
\]
\[
\rho_0 = \begin{pmatrix} 1 & -1 \\ 0 & 1 \end{pmatrix}, \qquad
\sigma(\varepsilon_0) = \rho(\varepsilon_0) = \varepsilon_1
\]
\note{$\tau''$ est écrit sur un $\tau_2$ biffé ; la lecture de l'indice de $\rho_0$, écrit sur une autre lettre, est douteuse.}

\section{Le groupe $D_\infty \times \mu$ agissant sur le cylindre}

\page{169}
\[
(z, s) \longrightarrow (\varepsilon z \exp 2i\pi a s, \ \varepsilon' s), \qquad \varepsilon, \varepsilon' \in \mu, \quad a \in \mathbb{Q}
\]
\drawing{un cylindre couché traversé par un axe vertical.}
\[
\bigl(z \exp 2i\pi a s\bigr)^{(-1)^\nu}, \quad \varepsilon s
\]
\struck{$\varepsilon \ \varepsilon'$} \quad $D_\infty \times \mu$

Dans un cadre : $\varepsilon = 1$ ; $\varepsilon = 1$, $\varepsilon' = -1$.

\drawing{un cylindre couché, traversé d'un axe, avec un segment oblique à l'intérieur.}
\[
\begin{aligned}
\sigma_1(z, s) &= (z^{-1}, s) \\
\sigma_2(z, s) &= (z, -s) \\
\theta(z, s) &= (z \exp 2i\pi s, s)
\end{aligned}
\qquad
\mathbf{a}(z, s) = (-z, -s)
\]
\note{la lettre $\mathbf{a}$ est soulignée sur la page ; elle est rendue ici en gras. Dans $\sigma_2(z,s)$, le signe devant $s$ surcharge un autre caractère ; devant $\theta$, une lettre biffée.}

Encadré :
\[
\begin{gathered}
\mathbf{a}^2 = \sigma_1^2 = \sigma_2^2 = 1, \qquad \sigma_1\sigma_2 = \sigma_2\sigma_1, \quad \mathbf{a}\sigma_1 = \sigma_1\mathbf{a} \\
\sigma_1(\theta) = \theta^{-1}, \qquad \sigma_2(\theta) = \theta^{-1} \\
\mathbf{a}(\theta) = \theta^{-1}
\end{gathered}
\]
\[
\mathbf{a}, \ \alpha, \ \beta, \ \theta
\]
\[
\begin{gathered}
(z^{-1} \exp 2i\pi s, \ s) \\
(z \exp -2i\pi s, \ s) \\
z \exp -2i\pi s, \ {+s} \\
-z \exp -2i\pi s, \ -s \\
z \exp -2i\pi s, \ s)
\end{gathered}
\]

\section{Composition dans une extension ; l'élément $\tilde\eta$}

\page{172}
\[
a \xrightarrow{\ l\ } \gamma^{-1}(a), \qquad a \xrightarrow{\ l'\ } \gamma'^{-1}(a)
\]
\[
(\gamma, l)(\gamma', l') = \bigl(\gamma\gamma', \ \gamma'^{-1}(l) \circ l'\bigr)
\]
\[
(\gamma . l)(\gamma' . l') = (\gamma\gamma')\,\underbrace{\gamma'^{-1}\, l\, \gamma'}\, l'
\]
\[
a \longrightarrow (\gamma\gamma')^{-1}a = \gamma'^{-1}(\gamma^{-1}a), \qquad
a \xrightarrow{\ l'\ } \gamma'^{-1}a \xrightarrow{\ \gamma'^{-1}(l)\ } (\gamma\gamma')^{-1}a
\]
\note{ces trois lignes sont au crayon, plus pâles que la suite.}
\[
\tilde u\,\tilde v\,\tilde u^{-1}\,\tilde v^{-1}
\]
\drawing{chemins $\gamma a \xrightarrow{l} a \to \gamma^{-1}a$ et $\gamma' a \xrightarrow{l'} a$ ; puis $\gamma\gamma' a \xrightarrow{\gamma(l')} \gamma(a) \xrightarrow{l} a$, et, biffé, $\gamma'(l') \to \gamma'(a)$, $e'$.}
\[
(\gamma, l)(\gamma', l') = \bigl(\gamma\gamma', \ l \circ \gamma(l')\bigr)
\]
\note{$l \circ \gamma(l')$ est écrit au-dessus d'une première expression biffée.}
\[
1 \longrightarrow \pi \longrightarrow E \longrightarrow G \longrightarrow 1
\]
\[
(l, \gamma)(l', \gamma') = (l . \gamma(l'))
\]
Encadré :
\[
(l, \gamma)(l', \gamma') = \bigl(l . \gamma(l'), \ \gamma\gamma'\bigr)
\]
\[
\text{\struck{$u \longmapsto u$, $a = \varepsilon_3(\tau')\alpha^{-1}u$, $b = \beta^{-1}u$}}
\qquad
\begin{aligned}
\alpha &= u\,\mathbf{a}^{-1}\varepsilon_3(\tau') \\
\beta &= u\,\mathbf{b}^{-1}
\end{aligned}
\]
\note{ici et plus bas, les lettres soulignées $\mathbf{a}$, $\mathbf{b}$, $\mathbf{u}$ sont rendues en gras. Le symbole $\mathfrak{z}$ transcrit une lettre en forme de 3 bouclé, qui note l'image dans $G$ d'un élément de $\widehat{\mathbb{Z}}$ (cf. p. 174).}
\[
\begin{aligned}
\tilde\eta &= \tilde\beta\,\tilde\sigma\,(\tilde\beta)^{-1} &&= \eta_0\,\mathfrak{z}(6\nu) \\
\tilde\sigma(\tilde\eta) &= \tilde\sigma(\tilde\beta)\,\tilde\beta^{-1} = \tilde\eta^{-1} &&= \tilde\sigma(\eta_0)\,\mathfrak{z}(6\nu) = \eta_0^{-1}\mathfrak{z}(6\nu)
\end{aligned}
\]
$\tilde\sigma^2$ ne commute à $\tilde\beta$ que si $\nu_2 = 0$, i.e. $\nu = 2\nu'$.
\[
\begin{aligned}
\tilde\eta = \eta_0 = [\beta_0, \sigma_0] &= \beta_0(\sigma_0)\sigma_0^{-1} = \mathbf{u}(\sigma_0)\sigma_0^{-1}, \qquad \tilde\sigma = \sigma_0\,\mathfrak{z}(3)
\end{aligned}
\]
\[
\text{\struck{$\tilde\eta = u(\tilde\sigma)\tilde\sigma^{-1} = u($}}
\]
\[
\begin{aligned}
\tilde\eta = \mathbf{u}(\tilde\sigma)\,\tilde\sigma^{-1}
&= \mathbf{u}(\sigma_0)\,\mathfrak{z}(3\mu)\,\sigma_0^{-1}\,\mathfrak{z}(-3) \\
&= \mathbf{u}(\sigma_0)\,\sigma_0^{-1}\,\mathfrak{z}\bigl(\underbrace{3(\mu - 1)}_{6\nu}\bigr)
\end{aligned}
\]
\[
\beta(\tilde\sigma) = u(\tilde\sigma), \qquad \beta(\tilde\sigma) = \beta(\sigma_0)\,\mathfrak{z}(3)
\]
\note{au bas de la page, sous $u(\tilde\sigma)$, un signe $=$ vertical suivi d'un trait inachevé.}

\page{173}
\[
\tilde\eta = \eta_0, \qquad \tilde\xi = \xi_0
\]
\[
\begin{aligned}
\tilde u(\tilde\varepsilon_0) = \tilde u(\tilde\sigma\tilde\rho)
&= \mathfrak{z}(\nu')\,\tilde\eta\,\tilde\sigma\,\mathfrak{z}(-\nu'')\,\tilde\xi\,\tilde\rho \\
&= \mathfrak{z}(\nu' - \nu'')\,\underbrace{\eta_0\,\sigma(\xi_0)}_{l_0^{\nu}}\,
\underbrace{\tilde\sigma\tilde\rho}_{\sigma_0\rho_0\mathfrak{z}(1) = \varepsilon_0\mathfrak{z}(1)} \\
&= \mathfrak{z}(\nu' - \nu'')\,\underbrace{l_0^{\nu}}_{\varepsilon_0^{2\nu}}\,\varepsilon_0\,\mathfrak{z}(1) \\
&= \mathfrak{z}(1 + \nu' - \nu'')\,\varepsilon_0^{2\nu+1} \\
&= \mathfrak{z}(1 + \nu' - \nu'')\,\underbrace{\varepsilon_0^{\mu}}_{\tilde\varepsilon_0^{\mu}\mathfrak{z}(-\mu)}
\end{aligned}
\]
\note{à la première ligne, un symbole biffé devant $\mathfrak{z}(\nu')$ ; à la quatrième, un facteur biffé devant $\varepsilon_0^{2\nu+1}$.}

Encadré :
\[
\tilde u(\tilde\varepsilon_0) = \mathfrak{z}(\nu' - \nu'' - 2\nu)\,\tilde\varepsilon_0^{\mu}
\]
\[
\begin{aligned}
&\underbrace{\nu' - \nu''}_{-\nu_2\nu_3 + \nu_0} = 2\nu \\
&-\nu_2\nu_3 + \nu_0 = 2(2\nu_3 - 3\nu_2 - 6\nu_2\nu_3 + 6\nu_0) \\
&11\nu_0 = 0, \qquad \nu_0 = 0
\end{aligned}
\]
\[
\begin{aligned}
\tilde u_1(\tilde\varepsilon_0) &= \mathfrak{z}(\nu_1' - \nu_1'' - 2\nu_1)\,\tilde\varepsilon_0^{\mu_1} \\
\tilde u_1\tilde u(\tilde\varepsilon_0) &= \mathfrak{z}\bigl(\mu_1(\nu' - \nu'' - 2\nu)\bigr)\,\tilde u_1(\tilde\varepsilon_0^{\mu})
\end{aligned}
\]
\note{le facteur $2\nu$ de l'encadré surcharge une première écriture. La ligne $11\nu_0 = 0$ ne se déduit pas telle quelle de la précédente ; elle est transcrite comme elle est écrite.}

\section{Construction}

\page{174}
\textbf{Construction.} Soit
\[
(1) \qquad G = Gl(2, \widehat{\mathbb{Z}}) . \widehat{\mathbb{Z}} \qquad \text{(produit semi-direct)}
\]
où $Gl(2, \widehat{\mathbb{Z}})$ opère sur $\widehat{\mathbb{Z}}$ par
\[
(2) \qquad u . \lambda = \det(u)\,\lambda
\]
\struck{Considérons \ill{}} Pour $u \in Gl(2, \widehat{\mathbb{Z}})$, on désigne par $u_G$ son image dans $G$, et pour $\lambda \in \widehat{\mathbb{Z}}$, on désigne par $\mathfrak{z}(\lambda)$ son image dans $G$.
\note{« produit ½ direct » sur la page. La lecture $u_G$ est douteuse.}
\[
(3) \qquad
\left\{
\begin{aligned}
&Sl(2, \mathbb{Z})^{\sim} \longrightarrow G \\
&\tilde\rho \longmapsto \rho_0 . \mathfrak{z}(-2) \\
&\tilde\sigma \longmapsto \sigma_0 . \mathfrak{z}(3) \\
&\Bigl[\ \tilde\omega_0 = \tilde\rho^{-3} = \tilde\sigma^2 \longmapsto \omega_0 . \mathfrak{z}(6), \quad
l'_0 = \tilde\omega_0^{2} \longmapsto \mathfrak{z}(12)\ \Bigr] \\
&\tilde r_0 = \tilde\tau' \longmapsto r_0
\end{aligned}
\right.
\]
\note{dans $\tilde\sigma \mapsto \sigma_0 . \mathfrak{z}(3)$, le $3$ surcharge un premier chiffre ; devant $r_0$, une lettre biffée.}

Soit le sous-groupe $SG$ \note{une première écriture $Sl(2,\mathbb{Z}) \ldots$ biffée}
\[
(4) \qquad SG = Sl(2, \widehat{\mathbb{Z}}) . \widehat{\mathbb{Z}}
\]
\uncertain{image par l'inverse} de $Sl(2, \widehat{\mathbb{Z}})$ dans $G$, \struck{opérant} \uncertain{admet comme} centre $\mu \times \widehat{\mathbb{Z}}$ --- il s'ensuit que l'opération de $G$ sur $SG$ par automorph. int. se factorise en une opération de $G/\widehat{\mathbb{Z}}^{\ast} \simeq Gl(2, \widehat{\mathbb{Z}})$ sur $SG$, donc si $\alpha \in Gl(2, \widehat{\mathbb{Z}})$, $x \in SG$, $\alpha(x) = \alpha x \alpha^{-1}$ est défini \struck{\ill{}}, \uncertain{à savoir} $\tilde\alpha x \tilde\alpha^{-1}$, où $\tilde\alpha$ \uncertain{est un} relèvement \uncertain{arbitraire} de $\alpha$ dans $G$ :
\[
(5) \qquad \alpha(x) \text{ ou } \alpha x \alpha^{-1} \overset{\text{déf}}{=} \tilde\alpha(x) = \tilde\alpha x \tilde\alpha^{-1} = \alpha_0(x) = \alpha_0 x \alpha_0^{-1}.
\]
Considérons
\[
(6) \qquad
\left\{
\begin{aligned}
&\beta \in Sl(2, \widehat{\mathbb{Z}}), \qquad \beta = \mathbf{u}\,b^{-1}, \quad \mathbf{u} = \begin{pmatrix} \mu & \lambda \\ 0 & 1 \end{pmatrix}, \quad b \in \mathrm{Cent}(L_\sigma) \\
&\text{i.e.}\quad \beta = \mathbf{u}\,(t\sigma + u)^{-1}
\end{aligned}
\right.
\]
\[
\left|
\begin{aligned}
&\nu_2 \in \{0, 1\}, \quad t^2 + u^2 = \mu \in \widehat{\mathbb{Z}}^{\ast} \\
&t, u \in \widehat{\mathbb{Z}}
\end{aligned}
\right.
\]
\note{après $Sl(2,\widehat{\mathbb{Z}})$, un mot raturé ; sous $\mathbf{u}$, une ligne biffée ; sous $\mathrm{Cent}(L_\sigma)$, la précision $(Gl(2,\widehat{\mathbb{Z}}))$. La lecture $\mathrm{Cent}$ est douteuse.}

Considérons un rel. $\tilde\beta$ de $\beta$ dans $SG$ (p. ex. $\tilde\beta = \beta_0$), d'où
\[
(7) \qquad \tilde\eta \overset{\text{déf}}{=} [\tilde\beta, \tilde\sigma] = \tilde\beta\,\tilde\sigma\,(\tilde\beta)^{-1}
= \tilde\beta(\tilde\sigma)\,\tilde\sigma^{-1} = \beta(\tilde\sigma)\,\tilde\sigma^{-1}
\]
\note{devant $[\tilde\beta, \tilde\sigma]$, une première écriture raturée.}

On a
\[
\sigma(\tilde\eta) = \tilde\sigma(\tilde\eta) = \tilde\sigma(\tilde\beta)\,\underbrace{\tilde\sigma^2(\tilde\beta)^{-1}} =
\]
\note{au-dessus du premier $=$, un « (déf) » biffé. La ligne s'arrête là, au bas du feuillet.}

\page{176}
Encadré :
\[
\left\{
\begin{aligned}
&\tilde u(\tilde\rho) = \lambda'\,\tilde\xi\,\tilde\rho \\
&\tilde u(\tilde\sigma) = \lambda''\,\tilde\eta\,\tilde\sigma \\
&\tilde u(\tilde\omega_0) = \tilde\omega_0^{\mu} \\
&\tilde u(\tilde\varepsilon_0) \overset{?}{=} \tilde\varepsilon_0^{\mu}
\end{aligned}
\right.
\quad\Longleftrightarrow\quad
\underbrace{\nu' - \nu''}_{-\nu_2\nu_3 + \nu_0} = 2\nu
\quad (\text{i.e. } \nu_0 = 0 \text{ si } \nu_2 = \nu_3 = 0)
\]
$\lambda', \lambda'' \in$ \uncertain{$\widehat{Z}_{1,1}$}.
\note{dans la deuxième ligne de l'encadré, $\lambda''$ surcharge une autre lettre ; sous l'accolade, une ligne biffée.}

\textbf{Ext} des gr. de $\mathcal{M}_{0,3}[0]$ dans $Sl(2, \mathbb{Z})^{\wedge}$ : $Sl(2, \mathbb{Z})^{\sim\wedge}$

Relations
\[
\left\{
\begin{aligned}
&\tilde\rho^{-3} = \tilde\sigma^2 = \tilde\omega_0 \\
&\tilde\varepsilon_0 = \tilde\sigma\tilde\rho
\end{aligned}
\right.
\qquad \text{N.B. on pose } l'_0 = \tilde\omega_0^2
\]
\struck{$\lambda'$, $\lambda''$ déterminés par}
\[
\text{\struck{$\underbrace{\tilde\eta\,\tilde\sigma(\tilde\eta)}_{\in \widehat{Z}_{1,1}}\,\lambda''^2 = l_0\,\tilde\omega_0^{2\nu}$,
\quad $\underbrace{\tilde\xi\,\rho(\tilde\xi)\,\rho^2(\tilde\xi)}_{\in Z_{1,1}}\,\lambda'^3 = \tilde\omega_0^{2\nu}$}}
\]
\note{ce passage est barré d'une grande croix ; un terme raturé après $\tilde\sigma(\tilde\eta)$.}
\[
\begin{aligned}
\tilde u(\tilde\rho) &= \lambda'\,\mathrm{int}(\tilde\alpha)(\tilde\rho) \\
\tilde u(\tilde\rho^3) &= \lambda'^3\,\mathrm{int}(\tilde\alpha)(\underbrace{\tilde\rho^3}_{\tilde\omega_0^{-1}})
= \lambda'^3\,\tilde\omega_0^{-\varepsilon_3} \overset{?}{=} \tilde u(\tilde\omega_0^{-1}) = \tilde\omega_0^{-\mu} \\
\lambda'^3 &= \tilde\omega_0^{-(\mu - \varepsilon_3)}
\end{aligned}
\]
\note{à la fin de la dernière ligne, $= \tilde\omega_0^{-2\nu}$ biffé.}

Encadré :
\[
\begin{aligned}
\lambda' &= \tilde\omega_0^{-\frac{\mu - \varepsilon_3}{3}} = \tilde\omega_0^{-2\nu''} = l'^{-\nu''}_0 \\
\lambda'' &= \tilde\omega_0^{\frac{\mu - \varepsilon_2}{2}} = \omega_0^{2\nu'} = l'^{\nu'}_0
\end{aligned}
\]
\[
\lambda'\lambda'' = \omega_0^{\frac{\mu - 3\varepsilon_2 + 2\varepsilon_3}{6}}
\]
\note{les égalités à droite de l'encadré sont serrées dans la marge et surchargées ; leur lecture est douteuse.}
\[
\text{\struck{$\tilde u(\tilde\sigma) = \lambda''\,\mathrm{int}(\tilde\beta)(\tilde\sigma)$}}
\]
\[
\text{\struck{$\tilde u(\tilde\sigma^2) = \lambda''^2\,\mathrm{int}(\tilde\beta)(\tilde\sigma^2) = \lambda''^2\,\tilde\omega_0^{\varepsilon_2}$}}
\;\text{\struck{$\overset{?}{=} \tilde u(\tilde\omega_0) = \tilde\omega_0^{\mu}$}}
\]
\[
\begin{aligned}
\tilde u(\tilde\varepsilon_0) &= \tilde u(\tilde\sigma)\,\tilde u(\tilde\rho)
= \lambda''\,\tilde\eta\,\tilde\sigma\,\lambda'\,\tilde\xi\,\tilde\rho
= \lambda''\,\tilde\eta\,\tilde\sigma(\lambda'\tilde\xi)\,\tilde\sigma\tilde\rho \\
&= \lambda''\,\tilde\eta . \sigma(\lambda'\tilde\xi)\,\tilde\varepsilon_0 \overset{?}{=} \tilde\varepsilon_0^{\mu}
\end{aligned}
\]
\[
\lambda''\,\tilde\eta\,\sigma(\lambda'\tilde\xi) \overset{?}{=} \tilde l_0^{\nu}
\]
\[
(\lambda'\lambda'')\,\underbrace{\sigma(\tilde\eta)\,\tilde\xi}_{\tilde\eta^{-1}\tilde\xi} \overset{?}{=} \tilde l_1^{\nu}
\qquad
\begin{aligned}
\tilde\eta &= \tilde\beta\,\sigma(\tilde\beta^{-1}) \\
\sigma(\tilde\eta) &= \sigma(\tilde\beta)\,\tilde\beta^{-1} = \tilde\eta^{-1}
\end{aligned}
\]
Encadré :
\[
(\lambda'\lambda'') \overset{?}{=} \tilde\xi^{-1}\,\tilde\eta\,\tilde l_1^{\nu} \quad ?
\]
\note{dans l'encadré, un terme raturé après $(\lambda'\lambda'')$ ; le $\tilde\xi$ est surchargé.}

\section{Les groupes $\Gamma^{!+}_{g,\nu}$ et leurs sous-groupes canoniques}

\page{179}
\[
\begin{array}{ccccccccc}
\mathfrak{S}_\nu & & \mathfrak{S}_{\nu-1} & & \mathfrak{S}_5 & & \mathfrak{S}_4 & & \\
\Gamma^{!+}_{0,\nu} & \longrightarrow & \Gamma^{!+}_{0,\nu-1} & \cdots & \Gamma^{!+}_{0,5} & \longrightarrow & \Gamma^{!+}_{0,4} & \longrightarrow & \Gamma^{!+}_{0,3} \\
\cup & & \cup & & \cup & & \wr & & \| \\
 & & & & & & Sl(2,\mathbb{Z})/\pm 1 & & \{1\} \\
 & & & & & & \cup & & \\
\Pi_{0,\nu-1} & & \Pi_{0,\nu-2} & \cdots & \Pi_{0,4} & & \Pi_{0,3} & & \\
\cup & & \cup & & \cup & & \cup & & \\
L_0^{(0,\nu-1)} & & L_0^{(0,\nu-1)} & & L_0^{(0,4)} & & L_0^{(0,3)} & &
\end{array}
\]
\note{les $\mathfrak{S}_\nu, \ldots, \mathfrak{S}_4$ sont écrits au-dessus des groupes. Sous $\Gamma^{!+}_{0,3}$, un premier groupe raturé, remplacé par $\{1\}$. Le deuxième exposant $(0, \nu - 1)$ est tel qu'on le lit ; on attendrait $(0, \nu - 2)$.}

Dans $\Gamma^{!+}_{0,\nu}$ \add{$(\nu \geq 4)$}, sous-groupes canoniques \struck{nilpotents} commutatifs \add{\uncertain{libres}} de rang $\nu - 3$ (i.e. $\simeq \mathbb{Z}^{\nu-3}$).

\drawing{une rangée de points numérotés $0, 1, 2, 3, \ldots, \nu - 1, \nu$ ; au-dessous, un cercle contenant deux points, puis deux points isolés marqués $1$ et $\infty$.}
\[
k[x_1, x_2, \ldots, x_\nu] \subset k\{x_1\}\{x_2\} \cdots \{x_\nu\}
\]
\note{un terme raturé entre $\{x_1\}$ et les points de suspension.}
\[
\begin{array}{ccccccc}
\Gamma^{!+}_{1,\nu} & \longrightarrow & \cdots & \Gamma^{!+}_{1,3} & \longrightarrow & \Gamma^{!+}_{1,2} \longrightarrow & \Gamma^{!+}_{1,1} \\
 & & & \cup & & \cup & \wr \\
\Pi_{1,\nu-1} & & & \Pi_{1,2} & & \Pi_{1,1} & Sl(2,\mathbb{Z}) \\
\cup & & & \cup & & \cup & \cup \\
L_0^{(1,\nu-1)} & & & L_0^{(1,2)} & & L_0^{(1,1)} & L_0^{\ill{}}
\end{array}
\]
Dans $\Gamma^{!+}_{1,\nu}$ $(\nu \geq 1)$, sous-groupes can. \struck{nilpotents} \ill{} \struck{\ill{}} de rang $\nu$.

$g \geq 2$, $\nu \geq 0$
\[
\begin{array}{cccccc}
\Gamma^{!+}_{g,\nu} \longrightarrow \cdots & & \Gamma^{!+}_{g,2} & \longrightarrow \Gamma^{!+}_{g,1} & \longrightarrow \Gamma^{!+}_{g,0} & \\
\cup & & \cup & \cup & \cup & \\
\Pi_{g,\nu-1} & \cdots & \Pi_{g,1} & \Pi_g & \Delta_g \simeq \mathbb{Z}^g & \\
\cup & & \cup & \cup & & \\
L_0^{(g,\nu-1)} & & L_0^{(g,1)} & \Lambda_g\ (\simeq \mathbb{Z}) & & \\
 & & & ? & &
\end{array}
\]
Sous-groupes canoniques \struck{nilpotents} \uncertain{maximaux} de rang \uncertain{$3g + \nu$}.
\note{la lettre $\Delta_g$ (tracée en gras) est d'une lecture douteuse. L'exposant de $\Gamma^{!+}_{g,\nu}$ est écrit « $!+1$ ». Le chiffre devant $g$ dans « rang $3g+\nu$ » surcharge un autre chiffre ; la lecture est douteuse. Au bas de la page, une formule entièrement biffée, illisible, se terminant par $\frac{(\nu - 1)g}{2}$ \uncertain{environ}.}

\end{document}
